% GENERATED FILE. Edit canonical lessons, then run npm run generate:books.
% canonical-source-hash: fnv1a64:1c9a34a5b3300683
% canonical-lessons: AR-C335-sahir, AR-C335-ghul, AR-C335-tinnin, AR-C335-imlaq, AR-C335-qazam

\chapter{Fairy-tale creatures}
\label{ch:ar-fairy-tale-creatures}
\hlchaptermodality{335}

\begin{chapteropening}
\textbf{By the end of this chapter:} \emph{I can say a magician, an ogre, a dragon, a giant and a dwarf.}

Complete the last lesson of chapter 335: I can say a magician, an ogre, a dragon, a giant and a dwarf.
\end{chapteropening}

\section[\texorpdfstring{\ar{ساحر} (sāḥir)}{ساحر (sāḥir)}]{\ar{ساحر} (sāḥir) — a magician, a sorcerer}

\label{lesson:AR-C335-sahir}



\begin{quote}
Before the new one: say the Arabic for to convict, then the Arabic for innocence.
\end{quote}

\subsection*{You'll want to know: \ar{ساحر}}
\textbf{\ar{ساحر}} — \emph{sāḥir} — \textquotedblleft{}a magician, a sorcerer\textquotedblright{}.

\begin{quote}
\ar{حول} \ar{الساحر} \ar{الأمير} \ar{إلى} \ar{ضفدع}.
\end{quote}

\begin{quote}
\emph{ḥawwala as-sāḥir al-amīr ilā ḍifdaʿ.}
\end{quote}

\begin{quote}
The magician turned the prince into a frog.
\end{quote}

Plural: \ar{سحرة}, \emph{saḥara}.

Of a place or a smile, \ar{ساحر} is enchanting.

\begin{cousinweb}
\textbf{\ar{ساحر}} (\emph{sāḥir}) is built on the native Arabic root \emph{s-ḥ-r} (Lane), \textquotedblleft{}to bewitch, to enchant\textquotedblright{}; \emph{siḥr} is magic and \emph{sāḥir} the one who works it.
\end{cousinweb}

\begin{grammarlens}[title={What you've built}]
One more word for telling what happened.
\end{grammarlens}

\subsection*{Your turn}
\begin{itemize}
\raggedright
  \item \emph{Say it:} \emph{sāḥir}
  \item \emph{Say it:} \emph{sāḥir}, once more
  \item \emph{Say it:} \emph{ḥawwala as-sāḥir al-amīr ilā ḍifdaʿ.} — The magician turned the prince into a frog.
  \item \emph{Recall:} say the Arabic for to convict, then the Arabic for innocence, then say \emph{sāḥir} again
\end{itemize}

\begin{tcolorbox}[breakable,colback=teal!4,colframe=teal!35!black,title={Before you move on}]
What does \ar{ساحر} mean? (\textquotedblleft{}A magician, a sorcerer\textquotedblright{}.) Say it once more.
\end{tcolorbox}
\section[\texorpdfstring{\ar{غول} (ghūl)}{غول (ghūl)}]{\ar{غول} (ghūl) — an ogre, a ghoul}

\label{lesson:AR-C335-ghul}



\begin{quote}
Before the new one: say the Arabic for innocence, then the Arabic for a magician.
\end{quote}

\subsection*{You'll want to know: \ar{غول}}
\textbf{\ar{غول}} — \emph{ghūl} — \textquotedblleft{}an ogre, a ghoul\textquotedblright{}.

\begin{quote}
\ar{خاف} \ar{الأطفال} \ar{من} \ar{الغول} \ar{في} \ar{الحكاية}.
\end{quote}

\begin{quote}
\emph{khāfa al-aṭfāl min al-ghūl fī al-ḥikāya.}
\end{quote}

\begin{quote}
The children were scared of the ogre in the tale.
\end{quote}

Plural: \ar{غيلان}, \emph{ghīlān}.

English ghoul is usually traced to this word.

\begin{cousinweb}
\textbf{\ar{غول}} (\emph{ghūl}) is built on the native Arabic root \emph{gh-w-l} (Lane), \textquotedblleft{}to take and destroy without warning\textquotedblright{}. Lane gives \emph{ghūl} as a being of the desert believed to lead travellers astray.
\end{cousinweb}

\begin{grammarlens}[title={What you've built}]
One more word for telling what happened.
\end{grammarlens}

\subsection*{Your turn}
\begin{itemize}
\raggedright
  \item \emph{Say it:} \emph{ghūl}
  \item \emph{Say it:} \emph{ghūl}, once more
  \item \emph{Say it:} \emph{khāfa al-aṭfāl min al-ghūl fī al-ḥikāya.} — The children were scared of the ogre in the tale.
  \item \emph{Recall:} say the Arabic for innocence, then the Arabic for a magician, then say \emph{ghūl} again
\end{itemize}

\begin{tcolorbox}[breakable,colback=teal!4,colframe=teal!35!black,title={Before you move on}]
What does \ar{غول} mean? (\textquotedblleft{}An ogre, a ghoul\textquotedblright{}.) Say it once more.
\end{tcolorbox}
\section[\texorpdfstring{\ar{تنين} (tinnīn)}{تنين (tinnīn)}]{\ar{تنين} (tinnīn) — a dragon}

\label{lesson:AR-C335-tinnin}



\begin{quote}
Before the new one: say the Arabic for a magician, then the Arabic for an ogre.
\end{quote}

\subsection*{You'll want to know: \ar{تنين}}
\textbf{\ar{تنين}} — \emph{tinnīn} — \textquotedblleft{}a dragon\textquotedblright{}.

\begin{quote}
\ar{طار} \ar{التنين} \ar{فوق} \ar{الجبل}.
\end{quote}

\begin{quote}
\emph{ṭāra at-tinnīn fawqa al-jabal.}
\end{quote}

\begin{quote}
The dragon flew over the mountain.
\end{quote}

\begin{cousinweb}
\textbf{\ar{تنين}} (\emph{tinnīn}) is given by Lane, under \emph{t-n-n}, as a huge serpent. It was borrowed from Aramaic \emph{tannīnā} (EtymDB); Hebrew has \emph{tannīn}, a sea monster.
\end{cousinweb}

\begin{grammarlens}[title={What you've built}]
One more word for telling what happened.
\end{grammarlens}

\subsection*{Your turn}
\begin{itemize}
\raggedright
  \item \emph{Say it:} \emph{tinnīn}
  \item \emph{Say it:} \emph{tinnīn}, once more
  \item \emph{Say it:} \emph{ṭāra at-tinnīn fawqa al-jabal.} — The dragon flew over the mountain.
  \item \emph{Recall:} say the Arabic for a magician, then the Arabic for an ogre, then say \emph{tinnīn} again
\end{itemize}

\begin{tcolorbox}[breakable,colback=teal!4,colframe=teal!35!black,title={Before you move on}]
What does \ar{تنين} mean? (\textquotedblleft{}A dragon\textquotedblright{}.) Say it once more.
\end{tcolorbox}
\section[\texorpdfstring{\ar{عملاق} (ʿimlāq)}{عملاق (ʿimlāq)}]{\ar{عملاق} (ʿimlāq) — a giant}

\label{lesson:AR-C335-imlaq}



\begin{quote}
Before the new one: say the Arabic for an ogre, then the Arabic for a dragon.
\end{quote}

\subsection*{You'll want to know: \ar{عملاق}}
\textbf{\ar{عملاق}} — \emph{ʿimlāq} — \textquotedblleft{}a giant\textquotedblright{}.

\begin{quote}
\ar{في} \ar{القصة} \ar{عملاق} \ar{يعيش} \ar{في} \ar{الغابة}.
\end{quote}

\begin{quote}
\emph{fī al-qiṣṣa ʿimlāq yaʿīshu fī al-ghāba.}
\end{quote}

\begin{quote}
In the story a giant lives in the forest.
\end{quote}

Plural: \ar{عمالقة}, \emph{ʿamāliqa}.

Of a company or a team: \ar{عملاق}, a giant.

\begin{cousinweb}
\textbf{\ar{عملاق}} (\emph{ʿimlāq}) is built on the four-letter root \emph{ʿ-m-l-q} (Lane), where it is \textquotedblleft{}very tall\textquotedblright{}, and the \emph{ʿamāliqa} are a people of old remembered as giants. Its origin before that is not settled.
\end{cousinweb}

\begin{grammarlens}[title={What you've built}]
One more word for telling what happened.
\end{grammarlens}

\subsection*{Your turn}
\begin{itemize}
\raggedright
  \item \emph{Say it:} \emph{ʿimlāq}
  \item \emph{Say it:} \emph{ʿimlāq}, once more
  \item \emph{Say it:} \emph{fī al-qiṣṣa ʿimlāq yaʿīshu fī al-ghāba.} — In the story a giant lives in the forest.
  \item \emph{Recall:} say the Arabic for an ogre, then the Arabic for a dragon, then say \emph{ʿimlāq} again
\end{itemize}

\begin{tcolorbox}[breakable,colback=teal!4,colframe=teal!35!black,title={Before you move on}]
What does \ar{عملاق} mean? (\textquotedblleft{}A giant\textquotedblright{}.) Say it once more.
\end{tcolorbox}
\section[\texorpdfstring{\ar{قزم} (qazam)}{قزم (qazam)}]{\ar{قزم} (qazam) — a dwarf}

\label{lesson:AR-C335-qazam}



\begin{quote}
Before the new one: say the Arabic for a dragon, then the Arabic for a giant.
\end{quote}

\subsection*{You'll want to know: \ar{قزم}}
\textbf{\ar{قزم}} — \emph{qazam} — \textquotedblleft{}a dwarf\textquotedblright{}.

\begin{quote}
\ar{ساعد} \ar{الأقزام} \ar{البنت} \ar{في} \ar{الغابة}.
\end{quote}

\begin{quote}
\emph{sāʿada al-aqzām al-bint fī al-ghāba.}
\end{quote}

\begin{quote}
The dwarfs helped the girl in the forest.
\end{quote}

Plural: \ar{أقزام}, \emph{aqzām}.

\begin{cousinweb}
\textbf{\ar{قزم}} (\emph{qazam}) is usually traced to a native Arabic root \emph{q-z-m}, where \emph{qazam} is something small and mean; Lane's article on it is empty. It now names a dwarf.
\end{cousinweb}

\begin{grammarlens}[title={What you've built}]
One more word for telling what happened.
\end{grammarlens}

\subsection*{Your turn}
\begin{itemize}
\raggedright
  \item \emph{Say it:} \emph{qazam}
  \item \emph{Say it:} \emph{qazam}, once more
  \item \emph{Say it:} \emph{sāʿada al-aqzām al-bint fī al-ghāba.} — The dwarfs helped the girl in the forest.
  \item \emph{Recall:} say the Arabic for a dragon, then the Arabic for a giant, then say \emph{qazam} again
\end{itemize}

\begin{tcolorbox}[breakable,colback=teal!4,colframe=teal!35!black,title={Before you move on}]
What does \ar{قزم} mean? (\textquotedblleft{}A dwarf\textquotedblright{}.) Say it once more.
\end{tcolorbox}
